% report-base.tex: preambolo comune per i rapporti PDF in stile nota Nomisma.
% Uso: nel file principale definire le macro di configurazione (\ReportLang,
% \ReportBrand, \ReportProject, \ReportTitle, \ReportSubtitle, \ReportAuthor,
% \ReportRole, \ReportDate), poi % report-base.tex: preambolo comune per i rapporti PDF in stile nota Nomisma.
% Uso: nel file principale definire le macro di configurazione (\ReportLang,
% \ReportBrand, \ReportProject, \ReportTitle, \ReportSubtitle, \ReportAuthor,
% \ReportRole, \ReportDate), poi % report-base.tex: preambolo comune per i rapporti PDF in stile nota Nomisma.
% Uso: nel file principale definire le macro di configurazione (\ReportLang,
% \ReportBrand, \ReportProject, \ReportTitle, \ReportSubtitle, \ReportAuthor,
% \ReportRole, \ReportDate), poi % report-base.tex: preambolo comune per i rapporti PDF in stile nota Nomisma.
% Uso: nel file principale definire le macro di configurazione (\ReportLang,
% \ReportBrand, \ReportProject, \ReportTitle, \ReportSubtitle, \ReportAuthor,
% \ReportRole, \ReportDate), poi \input{report-base} e \begin{document}.
% Compilare con xelatex, due passaggi (tre se cambiano riferimenti o bibliografia).
\documentclass[10pt,a4paper]{article}

% Valori di default: il file principale li sovrascrive definendoli prima di \input.
\providecommand{\ReportLang}{italian}% italian | english
\providecommand{\ReportBrand}{Nomisma}
\providecommand{\ReportProject}{}
\providecommand{\ReportTitle}{<<TITOLO>>}
\providecommand{\ReportSubtitle}{<<SOTTOTITOLO>>}
\providecommand{\ReportAuthor}{<<AUTORE>>}
\providecommand{\ReportRole}{<<RUOLO>>}
\providecommand{\ReportDate}{<<DATA>>}

\usepackage{etoolbox}
\usepackage{fontspec}
% Testo estraibile corretto dal PDF: senza questa riga le parentesi di Segoe UI
% vengono copiate come caratteri ornamentali (U+FD3E/U+FD3F).
\XeTeXgenerateactualtext=1

% Font: Segoe UI di Windows come nella nota di riferimento; fuori da Windows si
% ripiega su Latin Modern Sans, sempre presente in MiKTeX e TeX Live.
\IfFileExists{C:/Windows/Fonts/segoeui.ttf}{%
  \setmainfont{segoeui}[
    Path = C:/Windows/Fonts/,
    Extension = .ttf,
    BoldFont = segoeuib,
    ItalicFont = segoeuii,
    BoldItalicFont = segoeuiz]
  \setsansfont{segoeui}[
    Path = C:/Windows/Fonts/,
    Extension = .ttf,
    BoldFont = segoeuib,
    ItalicFont = segoeuii,
    BoldItalicFont = segoeuiz]
}{%
  \setmainfont{Latin Modern Sans}
  \setsansfont{Latin Modern Sans}
}

\usepackage[english,italian]{babel}
\AtBeginDocument{\selectlanguage{\ReportLang}}

\usepackage[a4paper,top=2.3cm,bottom=2.2cm,left=2.2cm,right=2.2cm,headheight=14pt,headsep=0.6cm,footskip=1cm]{geometry}
\usepackage[table]{xcolor}
\usepackage{colortbl}
\usepackage{array}
\usepackage{tabularx}
\usepackage{xltabular}
\usepackage{titlesec}
\usepackage{fancyhdr}
\usepackage{enumitem}
\usepackage{microtype}
\usepackage{needspace}
\usepackage{titletoc}
\usepackage[section]{placeins}
\usepackage{amsmath}
\usepackage{amssymb}
\usepackage{graphicx}
\usepackage{caption}
\usepackage{siunitx}

% Colori della nota metodologica
\definecolor{navy}{HTML}{0E2D53}
\definecolor{maroon}{HTML}{A3274B}
\definecolor{rowgrey}{HTML}{F2F2F2}
\definecolor{rulegrey}{HTML}{BFBFBF}
\definecolor{footgrey}{HTML}{8C8C8C}
\definecolor{boxgrey}{HTML}{F4F5F8}

\setlength{\parindent}{0pt}
\setlength{\parskip}{0.55em}
\linespread{1.12}

% Numeri: in italiano punto per le migliaia e virgola per i decimali, in inglese
% il contrario. Nel sorgente si scrive sempre il numero grezzo: \num{161250},
% \num{0.698}, \pct{96.6}. Il raggruppamento parte da 4 cifre (1.626).
% Modo testo: le cifre usano il font del testo e ereditano il grassetto.
% Incertezza esplicita: \num{0.874 +- 0.006} produce 0,874 ± 0,006.
\sisetup{mode=text,reset-text-family=false,reset-text-series=false,reset-text-shape=false,uncertainty-mode=separate}
\ifdefstring{\ReportLang}{english}{%
  \sisetup{output-decimal-marker={.}, group-separator={,}, group-minimum-digits=4}
  \newcommand{\ReportPageLabel}{p.}
}{%
  \sisetup{output-decimal-marker={,}, group-separator={.}, group-minimum-digits=4}
  \newcommand{\ReportPageLabel}{pag.}
}
\newcommand{\pct}[1]{\num{#1}\%}

% Testatina bordeaux in maiuscolo a sinistra, piede grigio a destra
\providecommand{\ReportHeader}{\ReportBrand\ifdefempty{\ReportProject}{}{\ \ | \ \ReportProject}}
\providecommand{\ReportFooter}{\ReportBrand\ifdefempty{\ReportProject}{}{ | \ReportProject}}
\newcommand{\ReportHeadStyle}{\footnotesize\bfseries\color{maroon}\MakeUppercase{\ReportHeader}}
\newcommand{\ReportFootStyle}{\scriptsize\color{footgrey}\ReportFooter\ ·\ \ReportPageLabel\ \thepage}
\pagestyle{fancy}
\fancyhf{}
\renewcommand{\headrulewidth}{0pt}
\fancyhead[L]{\ReportHeadStyle}
\fancyfoot[R]{\ReportFootStyle}
\fancypagestyle{plain}{\fancyhf{}\renewcommand{\headrulewidth}{0pt}%
  \fancyhead[L]{\ReportHeadStyle}\fancyfoot[R]{\ReportFootStyle}}

% Blocco del titolo: titolo blu in grassetto, sottotitolo, autore e ruolo, data
\newcommand{\maketitleblock}{%
  {\color{navy}\bfseries\LARGE \ReportTitle\par}
  \vspace{0.35em}
  {\color{navy}\large \ReportSubtitle\par}
  \vspace{0.5em}
  {\small \ReportAuthor, \ReportRole\par
  \color{footgrey}\ReportDate\par}}

% Titoli di sezione: numero e titolo blu in grassetto con filetto bordeaux sotto.
% Le sottosezioni sono blu, in grassetto, di corpo normale (piu' piccole delle sezioni).
\newcommand{\sectitle}[1]{\parbox[b]{\linewidth}{\thesection.\hspace{0.4em}#1\par\nointerlineskip\vspace{3pt}{\color{maroon}\rule{\linewidth}{0.7pt}}}}
\newcommand{\sectitlenn}[1]{\parbox[b]{\linewidth}{#1\par\nointerlineskip\vspace{3pt}{\color{maroon}\rule{\linewidth}{0.7pt}}}}
\titleformat{\section}[block]{\color{navy}\bfseries\large}{}{0pt}{\sectitle}
\titleformat{name=\section,numberless}[block]{\color{navy}\bfseries\large}{}{0pt}{\sectitlenn}
\titlespacing*{\section}{0pt}{1.0em}{0.2em}
\titleformat{\subsection}{\color{navy}\bfseries\normalsize}{\thesubsection}{0.6em}{}
\titlespacing*{\subsection}{0pt}{0.9em}{0.15em}
% Indice: sezioni in blu scuro, sottosezioni con puntini
\titlecontents{section}[1.6em]{\addvspace{0.4em}\color{navy}\bfseries}{\contentslabel{1.6em}}{\hspace*{-1.6em}}{\hfill\contentspage}
\titlecontents{subsection}[4.0em]{}{\contentslabel{2.4em}}{\hspace*{-2.4em}}{\titlerule*[0.7pc]{.}\contentspage}
% Evita titoli orfani in fondo pagina: se restano poche righe si va a capo pagina
\pretocmd{\section}{\needspace{9\baselineskip}}{}{}
\pretocmd{\subsection}{\needspace{5\baselineskip}}{}{}
% Le tabelle restano nella sezione in cui sono citate: placeins[section] chiude i
% float a ogni \section; a fine sezione si aggiunge comunque \FloatBarrier esplicito.
\renewcommand{\topfraction}{0.9}
\renewcommand{\bottomfraction}{0.8}
\renewcommand{\textfraction}{0.07}
\renewcommand{\floatpagefraction}{0.75}
\setcounter{topnumber}{3}
\setcounter{bottomnumber}{2}
\setcounter{totalnumber}{5}

\setlist{itemsep=0.15em,topsep=0.2em,parsep=0pt,leftmargin=1.3em}

% Didascalie delle tabelle: sopra la tabella, "Tabella N." in blu grassetto
\DeclareCaptionFont{navy}{\color{navy}}
\captionsetup{font=small,labelfont={bf,navy},textfont=navy,labelsep=period,justification=raggedright,singlelinecheck=false,skip=4pt,position=top}
\setlength{\LTcapwidth}{\linewidth}

% Tabelle: intestazione blu con testo bianco, righe alternate grigie, filetti sottili.
% Uso: \begin{tabella} ... tabularx o xltabular ... \end{tabella};
% riga di intestazione: \rowcolor{navy}\hd{Colonna} & ... \\ \hline
% xltabular va a capo pagina e ripete l'intestazione posta prima di \endhead.
\arrayrulecolor{rulegrey}
\setlength{\arrayrulewidth}{0.4pt}
\setlength{\LTpre}{0.3em}
\setlength{\LTpost}{0.3em}
\newcommand{\hd}[1]{\textcolor{white}{\bfseries #1}}
\newcolumntype{L}{>{\raggedright\arraybackslash}X}
\newcolumntype{P}[1]{>{\raggedright\arraybackslash}p{#1}}
\newcolumntype{R}[1]{>{\raggedleft\arraybackslash}p{#1}}
\newcolumntype{C}[1]{>{\centering\arraybackslash}p{#1}}
\newenvironment{tabella}{\par\footnotesize\renewcommand{\arraystretch}{1.3}\rowcolors{2}{rowgrey}{white}}{\par\normalsize}

% Riquadro di sintesi: bordo blu, fondo grigio chiaro (non si spezza tra pagine)
\newsavebox{\riquadrobox}
\newenvironment{riquadro}
  {\par\noindent\begin{lrbox}{\riquadrobox}\begin{minipage}{\dimexpr\linewidth-2\fboxsep-2\fboxrule\relax}\vspace{0.3em}}
  {\vspace{0.2em}\end{minipage}\end{lrbox}\fcolorbox{navy}{boxgrey}{\usebox{\riquadrobox}}\par}

% Collegamenti: hyperref per ultimo
\usepackage{hyperref}
\hypersetup{colorlinks=true,linkcolor=navy,citecolor=navy,urlcolor=maroon}
\AtBeginDocument{\hypersetup{pdftitle={\ReportTitle},pdfauthor={\ReportAuthor}}}
% cleveref dopo hyperref: \cref{tab:x} produce "Tabella 1"
\usepackage[capitalise,noabbrev]{cleveref}
 e \begin{document}.
% Compilare con xelatex, due passaggi (tre se cambiano riferimenti o bibliografia).
\documentclass[10pt,a4paper]{article}

% Valori di default: il file principale li sovrascrive definendoli prima di \input.
\providecommand{\ReportLang}{italian}% italian | english
\providecommand{\ReportBrand}{Nomisma}
\providecommand{\ReportProject}{}
\providecommand{\ReportTitle}{<<TITOLO>>}
\providecommand{\ReportSubtitle}{<<SOTTOTITOLO>>}
\providecommand{\ReportAuthor}{<<AUTORE>>}
\providecommand{\ReportRole}{<<RUOLO>>}
\providecommand{\ReportDate}{<<DATA>>}

\usepackage{etoolbox}
\usepackage{fontspec}
% Testo estraibile corretto dal PDF: senza questa riga le parentesi di Segoe UI
% vengono copiate come caratteri ornamentali (U+FD3E/U+FD3F).
\XeTeXgenerateactualtext=1

% Font: Segoe UI di Windows come nella nota di riferimento; fuori da Windows si
% ripiega su Latin Modern Sans, sempre presente in MiKTeX e TeX Live.
\IfFileExists{C:/Windows/Fonts/segoeui.ttf}{%
  \setmainfont{segoeui}[
    Path = C:/Windows/Fonts/,
    Extension = .ttf,
    BoldFont = segoeuib,
    ItalicFont = segoeuii,
    BoldItalicFont = segoeuiz]
  \setsansfont{segoeui}[
    Path = C:/Windows/Fonts/,
    Extension = .ttf,
    BoldFont = segoeuib,
    ItalicFont = segoeuii,
    BoldItalicFont = segoeuiz]
}{%
  \setmainfont{Latin Modern Sans}
  \setsansfont{Latin Modern Sans}
}

\usepackage[english,italian]{babel}
\AtBeginDocument{\selectlanguage{\ReportLang}}

\usepackage[a4paper,top=2.3cm,bottom=2.2cm,left=2.2cm,right=2.2cm,headheight=14pt,headsep=0.6cm,footskip=1cm]{geometry}
\usepackage[table]{xcolor}
\usepackage{colortbl}
\usepackage{array}
\usepackage{tabularx}
\usepackage{xltabular}
\usepackage{titlesec}
\usepackage{fancyhdr}
\usepackage{enumitem}
\usepackage{microtype}
\usepackage{needspace}
\usepackage{titletoc}
\usepackage[section]{placeins}
\usepackage{amsmath}
\usepackage{amssymb}
\usepackage{graphicx}
\usepackage{caption}
\usepackage{siunitx}

% Colori della nota metodologica
\definecolor{navy}{HTML}{0E2D53}
\definecolor{maroon}{HTML}{A3274B}
\definecolor{rowgrey}{HTML}{F2F2F2}
\definecolor{rulegrey}{HTML}{BFBFBF}
\definecolor{footgrey}{HTML}{8C8C8C}
\definecolor{boxgrey}{HTML}{F4F5F8}

\setlength{\parindent}{0pt}
\setlength{\parskip}{0.55em}
\linespread{1.12}

% Numeri: in italiano punto per le migliaia e virgola per i decimali, in inglese
% il contrario. Nel sorgente si scrive sempre il numero grezzo: \num{161250},
% \num{0.698}, \pct{96.6}. Il raggruppamento parte da 4 cifre (1.626).
% Modo testo: le cifre usano il font del testo e ereditano il grassetto.
% Incertezza esplicita: \num{0.874 +- 0.006} produce 0,874 ± 0,006.
\sisetup{mode=text,reset-text-family=false,reset-text-series=false,reset-text-shape=false,uncertainty-mode=separate}
\ifdefstring{\ReportLang}{english}{%
  \sisetup{output-decimal-marker={.}, group-separator={,}, group-minimum-digits=4}
  \newcommand{\ReportPageLabel}{p.}
}{%
  \sisetup{output-decimal-marker={,}, group-separator={.}, group-minimum-digits=4}
  \newcommand{\ReportPageLabel}{pag.}
}
\newcommand{\pct}[1]{\num{#1}\%}

% Testatina bordeaux in maiuscolo a sinistra, piede grigio a destra
\providecommand{\ReportHeader}{\ReportBrand\ifdefempty{\ReportProject}{}{\ \ | \ \ReportProject}}
\providecommand{\ReportFooter}{\ReportBrand\ifdefempty{\ReportProject}{}{ | \ReportProject}}
\newcommand{\ReportHeadStyle}{\footnotesize\bfseries\color{maroon}\MakeUppercase{\ReportHeader}}
\newcommand{\ReportFootStyle}{\scriptsize\color{footgrey}\ReportFooter\ ·\ \ReportPageLabel\ \thepage}
\pagestyle{fancy}
\fancyhf{}
\renewcommand{\headrulewidth}{0pt}
\fancyhead[L]{\ReportHeadStyle}
\fancyfoot[R]{\ReportFootStyle}
\fancypagestyle{plain}{\fancyhf{}\renewcommand{\headrulewidth}{0pt}%
  \fancyhead[L]{\ReportHeadStyle}\fancyfoot[R]{\ReportFootStyle}}

% Blocco del titolo: titolo blu in grassetto, sottotitolo, autore e ruolo, data
\newcommand{\maketitleblock}{%
  {\color{navy}\bfseries\LARGE \ReportTitle\par}
  \vspace{0.35em}
  {\color{navy}\large \ReportSubtitle\par}
  \vspace{0.5em}
  {\small \ReportAuthor, \ReportRole\par
  \color{footgrey}\ReportDate\par}}

% Titoli di sezione: numero e titolo blu in grassetto con filetto bordeaux sotto.
% Le sottosezioni sono blu, in grassetto, di corpo normale (piu' piccole delle sezioni).
\newcommand{\sectitle}[1]{\parbox[b]{\linewidth}{\thesection.\hspace{0.4em}#1\par\nointerlineskip\vspace{3pt}{\color{maroon}\rule{\linewidth}{0.7pt}}}}
\newcommand{\sectitlenn}[1]{\parbox[b]{\linewidth}{#1\par\nointerlineskip\vspace{3pt}{\color{maroon}\rule{\linewidth}{0.7pt}}}}
\titleformat{\section}[block]{\color{navy}\bfseries\large}{}{0pt}{\sectitle}
\titleformat{name=\section,numberless}[block]{\color{navy}\bfseries\large}{}{0pt}{\sectitlenn}
\titlespacing*{\section}{0pt}{1.0em}{0.2em}
\titleformat{\subsection}{\color{navy}\bfseries\normalsize}{\thesubsection}{0.6em}{}
\titlespacing*{\subsection}{0pt}{0.9em}{0.15em}
% Indice: sezioni in blu scuro, sottosezioni con puntini
\titlecontents{section}[1.6em]{\addvspace{0.4em}\color{navy}\bfseries}{\contentslabel{1.6em}}{\hspace*{-1.6em}}{\hfill\contentspage}
\titlecontents{subsection}[4.0em]{}{\contentslabel{2.4em}}{\hspace*{-2.4em}}{\titlerule*[0.7pc]{.}\contentspage}
% Evita titoli orfani in fondo pagina: se restano poche righe si va a capo pagina
\pretocmd{\section}{\needspace{9\baselineskip}}{}{}
\pretocmd{\subsection}{\needspace{5\baselineskip}}{}{}
% Le tabelle restano nella sezione in cui sono citate: placeins[section] chiude i
% float a ogni \section; a fine sezione si aggiunge comunque \FloatBarrier esplicito.
\renewcommand{\topfraction}{0.9}
\renewcommand{\bottomfraction}{0.8}
\renewcommand{\textfraction}{0.07}
\renewcommand{\floatpagefraction}{0.75}
\setcounter{topnumber}{3}
\setcounter{bottomnumber}{2}
\setcounter{totalnumber}{5}

\setlist{itemsep=0.15em,topsep=0.2em,parsep=0pt,leftmargin=1.3em}

% Didascalie delle tabelle: sopra la tabella, "Tabella N." in blu grassetto
\DeclareCaptionFont{navy}{\color{navy}}
\captionsetup{font=small,labelfont={bf,navy},textfont=navy,labelsep=period,justification=raggedright,singlelinecheck=false,skip=4pt,position=top}
\setlength{\LTcapwidth}{\linewidth}

% Tabelle: intestazione blu con testo bianco, righe alternate grigie, filetti sottili.
% Uso: \begin{tabella} ... tabularx o xltabular ... \end{tabella};
% riga di intestazione: \rowcolor{navy}\hd{Colonna} & ... \\ \hline
% xltabular va a capo pagina e ripete l'intestazione posta prima di \endhead.
\arrayrulecolor{rulegrey}
\setlength{\arrayrulewidth}{0.4pt}
\setlength{\LTpre}{0.3em}
\setlength{\LTpost}{0.3em}
\newcommand{\hd}[1]{\textcolor{white}{\bfseries #1}}
\newcolumntype{L}{>{\raggedright\arraybackslash}X}
\newcolumntype{P}[1]{>{\raggedright\arraybackslash}p{#1}}
\newcolumntype{R}[1]{>{\raggedleft\arraybackslash}p{#1}}
\newcolumntype{C}[1]{>{\centering\arraybackslash}p{#1}}
\newenvironment{tabella}{\par\footnotesize\renewcommand{\arraystretch}{1.3}\rowcolors{2}{rowgrey}{white}}{\par\normalsize}

% Riquadro di sintesi: bordo blu, fondo grigio chiaro (non si spezza tra pagine)
\newsavebox{\riquadrobox}
\newenvironment{riquadro}
  {\par\noindent\begin{lrbox}{\riquadrobox}\begin{minipage}{\dimexpr\linewidth-2\fboxsep-2\fboxrule\relax}\vspace{0.3em}}
  {\vspace{0.2em}\end{minipage}\end{lrbox}\fcolorbox{navy}{boxgrey}{\usebox{\riquadrobox}}\par}

% Collegamenti: hyperref per ultimo
\usepackage{hyperref}
\hypersetup{colorlinks=true,linkcolor=navy,citecolor=navy,urlcolor=maroon}
\AtBeginDocument{\hypersetup{pdftitle={\ReportTitle},pdfauthor={\ReportAuthor}}}
% cleveref dopo hyperref: \cref{tab:x} produce "Tabella 1"
\usepackage[capitalise,noabbrev]{cleveref}
 e \begin{document}.
% Compilare con xelatex, due passaggi (tre se cambiano riferimenti o bibliografia).
\documentclass[10pt,a4paper]{article}

% Valori di default: il file principale li sovrascrive definendoli prima di \input.
\providecommand{\ReportLang}{italian}% italian | english
\providecommand{\ReportBrand}{Nomisma}
\providecommand{\ReportProject}{}
\providecommand{\ReportTitle}{<<TITOLO>>}
\providecommand{\ReportSubtitle}{<<SOTTOTITOLO>>}
\providecommand{\ReportAuthor}{<<AUTORE>>}
\providecommand{\ReportRole}{<<RUOLO>>}
\providecommand{\ReportDate}{<<DATA>>}

\usepackage{etoolbox}
\usepackage{fontspec}
% Testo estraibile corretto dal PDF: senza questa riga le parentesi di Segoe UI
% vengono copiate come caratteri ornamentali (U+FD3E/U+FD3F).
\XeTeXgenerateactualtext=1

% Font: Segoe UI di Windows come nella nota di riferimento; fuori da Windows si
% ripiega su Latin Modern Sans, sempre presente in MiKTeX e TeX Live.
\IfFileExists{C:/Windows/Fonts/segoeui.ttf}{%
  \setmainfont{segoeui}[
    Path = C:/Windows/Fonts/,
    Extension = .ttf,
    BoldFont = segoeuib,
    ItalicFont = segoeuii,
    BoldItalicFont = segoeuiz]
  \setsansfont{segoeui}[
    Path = C:/Windows/Fonts/,
    Extension = .ttf,
    BoldFont = segoeuib,
    ItalicFont = segoeuii,
    BoldItalicFont = segoeuiz]
}{%
  \setmainfont{Latin Modern Sans}
  \setsansfont{Latin Modern Sans}
}

\usepackage[english,italian]{babel}
\AtBeginDocument{\selectlanguage{\ReportLang}}

\usepackage[a4paper,top=2.3cm,bottom=2.2cm,left=2.2cm,right=2.2cm,headheight=14pt,headsep=0.6cm,footskip=1cm]{geometry}
\usepackage[table]{xcolor}
\usepackage{colortbl}
\usepackage{array}
\usepackage{tabularx}
\usepackage{xltabular}
\usepackage{titlesec}
\usepackage{fancyhdr}
\usepackage{enumitem}
\usepackage{microtype}
\usepackage{needspace}
\usepackage{titletoc}
\usepackage[section]{placeins}
\usepackage{amsmath}
\usepackage{amssymb}
\usepackage{graphicx}
\usepackage{caption}
\usepackage{siunitx}

% Colori della nota metodologica
\definecolor{navy}{HTML}{0E2D53}
\definecolor{maroon}{HTML}{A3274B}
\definecolor{rowgrey}{HTML}{F2F2F2}
\definecolor{rulegrey}{HTML}{BFBFBF}
\definecolor{footgrey}{HTML}{8C8C8C}
\definecolor{boxgrey}{HTML}{F4F5F8}

\setlength{\parindent}{0pt}
\setlength{\parskip}{0.55em}
\linespread{1.12}

% Numeri: in italiano punto per le migliaia e virgola per i decimali, in inglese
% il contrario. Nel sorgente si scrive sempre il numero grezzo: \num{161250},
% \num{0.698}, \pct{96.6}. Il raggruppamento parte da 4 cifre (1.626).
% Modo testo: le cifre usano il font del testo e ereditano il grassetto.
% Incertezza esplicita: \num{0.874 +- 0.006} produce 0,874 ± 0,006.
\sisetup{mode=text,reset-text-family=false,reset-text-series=false,reset-text-shape=false,uncertainty-mode=separate}
\ifdefstring{\ReportLang}{english}{%
  \sisetup{output-decimal-marker={.}, group-separator={,}, group-minimum-digits=4}
  \newcommand{\ReportPageLabel}{p.}
}{%
  \sisetup{output-decimal-marker={,}, group-separator={.}, group-minimum-digits=4}
  \newcommand{\ReportPageLabel}{pag.}
}
\newcommand{\pct}[1]{\num{#1}\%}

% Testatina bordeaux in maiuscolo a sinistra, piede grigio a destra
\providecommand{\ReportHeader}{\ReportBrand\ifdefempty{\ReportProject}{}{\ \ | \ \ReportProject}}
\providecommand{\ReportFooter}{\ReportBrand\ifdefempty{\ReportProject}{}{ | \ReportProject}}
\newcommand{\ReportHeadStyle}{\footnotesize\bfseries\color{maroon}\MakeUppercase{\ReportHeader}}
\newcommand{\ReportFootStyle}{\scriptsize\color{footgrey}\ReportFooter\ ·\ \ReportPageLabel\ \thepage}
\pagestyle{fancy}
\fancyhf{}
\renewcommand{\headrulewidth}{0pt}
\fancyhead[L]{\ReportHeadStyle}
\fancyfoot[R]{\ReportFootStyle}
\fancypagestyle{plain}{\fancyhf{}\renewcommand{\headrulewidth}{0pt}%
  \fancyhead[L]{\ReportHeadStyle}\fancyfoot[R]{\ReportFootStyle}}

% Blocco del titolo: titolo blu in grassetto, sottotitolo, autore e ruolo, data
\newcommand{\maketitleblock}{%
  {\color{navy}\bfseries\LARGE \ReportTitle\par}
  \vspace{0.35em}
  {\color{navy}\large \ReportSubtitle\par}
  \vspace{0.5em}
  {\small \ReportAuthor, \ReportRole\par
  \color{footgrey}\ReportDate\par}}

% Titoli di sezione: numero e titolo blu in grassetto con filetto bordeaux sotto.
% Le sottosezioni sono blu, in grassetto, di corpo normale (piu' piccole delle sezioni).
\newcommand{\sectitle}[1]{\parbox[b]{\linewidth}{\thesection.\hspace{0.4em}#1\par\nointerlineskip\vspace{3pt}{\color{maroon}\rule{\linewidth}{0.7pt}}}}
\newcommand{\sectitlenn}[1]{\parbox[b]{\linewidth}{#1\par\nointerlineskip\vspace{3pt}{\color{maroon}\rule{\linewidth}{0.7pt}}}}
\titleformat{\section}[block]{\color{navy}\bfseries\large}{}{0pt}{\sectitle}
\titleformat{name=\section,numberless}[block]{\color{navy}\bfseries\large}{}{0pt}{\sectitlenn}
\titlespacing*{\section}{0pt}{1.0em}{0.2em}
\titleformat{\subsection}{\color{navy}\bfseries\normalsize}{\thesubsection}{0.6em}{}
\titlespacing*{\subsection}{0pt}{0.9em}{0.15em}
% Indice: sezioni in blu scuro, sottosezioni con puntini
\titlecontents{section}[1.6em]{\addvspace{0.4em}\color{navy}\bfseries}{\contentslabel{1.6em}}{\hspace*{-1.6em}}{\hfill\contentspage}
\titlecontents{subsection}[4.0em]{}{\contentslabel{2.4em}}{\hspace*{-2.4em}}{\titlerule*[0.7pc]{.}\contentspage}
% Evita titoli orfani in fondo pagina: se restano poche righe si va a capo pagina
\pretocmd{\section}{\needspace{9\baselineskip}}{}{}
\pretocmd{\subsection}{\needspace{5\baselineskip}}{}{}
% Le tabelle restano nella sezione in cui sono citate: placeins[section] chiude i
% float a ogni \section; a fine sezione si aggiunge comunque \FloatBarrier esplicito.
\renewcommand{\topfraction}{0.9}
\renewcommand{\bottomfraction}{0.8}
\renewcommand{\textfraction}{0.07}
\renewcommand{\floatpagefraction}{0.75}
\setcounter{topnumber}{3}
\setcounter{bottomnumber}{2}
\setcounter{totalnumber}{5}

\setlist{itemsep=0.15em,topsep=0.2em,parsep=0pt,leftmargin=1.3em}

% Didascalie delle tabelle: sopra la tabella, "Tabella N." in blu grassetto
\DeclareCaptionFont{navy}{\color{navy}}
\captionsetup{font=small,labelfont={bf,navy},textfont=navy,labelsep=period,justification=raggedright,singlelinecheck=false,skip=4pt,position=top}
\setlength{\LTcapwidth}{\linewidth}

% Tabelle: intestazione blu con testo bianco, righe alternate grigie, filetti sottili.
% Uso: \begin{tabella} ... tabularx o xltabular ... \end{tabella};
% riga di intestazione: \rowcolor{navy}\hd{Colonna} & ... \\ \hline
% xltabular va a capo pagina e ripete l'intestazione posta prima di \endhead.
\arrayrulecolor{rulegrey}
\setlength{\arrayrulewidth}{0.4pt}
\setlength{\LTpre}{0.3em}
\setlength{\LTpost}{0.3em}
\newcommand{\hd}[1]{\textcolor{white}{\bfseries #1}}
\newcolumntype{L}{>{\raggedright\arraybackslash}X}
\newcolumntype{P}[1]{>{\raggedright\arraybackslash}p{#1}}
\newcolumntype{R}[1]{>{\raggedleft\arraybackslash}p{#1}}
\newcolumntype{C}[1]{>{\centering\arraybackslash}p{#1}}
\newenvironment{tabella}{\par\footnotesize\renewcommand{\arraystretch}{1.3}\rowcolors{2}{rowgrey}{white}}{\par\normalsize}

% Riquadro di sintesi: bordo blu, fondo grigio chiaro (non si spezza tra pagine)
\newsavebox{\riquadrobox}
\newenvironment{riquadro}
  {\par\noindent\begin{lrbox}{\riquadrobox}\begin{minipage}{\dimexpr\linewidth-2\fboxsep-2\fboxrule\relax}\vspace{0.3em}}
  {\vspace{0.2em}\end{minipage}\end{lrbox}\fcolorbox{navy}{boxgrey}{\usebox{\riquadrobox}}\par}

% Collegamenti: hyperref per ultimo
\usepackage{hyperref}
\hypersetup{colorlinks=true,linkcolor=navy,citecolor=navy,urlcolor=maroon}
\AtBeginDocument{\hypersetup{pdftitle={\ReportTitle},pdfauthor={\ReportAuthor}}}
% cleveref dopo hyperref: \cref{tab:x} produce "Tabella 1"
\usepackage[capitalise,noabbrev]{cleveref}
 e \begin{document}.
% Compilare con xelatex, due passaggi (tre se cambiano riferimenti o bibliografia).
\documentclass[10pt,a4paper]{article}

% Valori di default: il file principale li sovrascrive definendoli prima di \input.
\providecommand{\ReportLang}{italian}% italian | english
\providecommand{\ReportBrand}{Nomisma}
\providecommand{\ReportProject}{}
\providecommand{\ReportTitle}{<<TITOLO>>}
\providecommand{\ReportSubtitle}{<<SOTTOTITOLO>>}
\providecommand{\ReportAuthor}{<<AUTORE>>}
\providecommand{\ReportRole}{<<RUOLO>>}
\providecommand{\ReportDate}{<<DATA>>}

\usepackage{etoolbox}
\usepackage{fontspec}
% Testo estraibile corretto dal PDF: senza questa riga le parentesi di Segoe UI
% vengono copiate come caratteri ornamentali (U+FD3E/U+FD3F).
\XeTeXgenerateactualtext=1

% Font: Segoe UI di Windows come nella nota di riferimento; fuori da Windows si
% ripiega su Latin Modern Sans, sempre presente in MiKTeX e TeX Live.
\IfFileExists{C:/Windows/Fonts/segoeui.ttf}{%
  \setmainfont{segoeui}[
    Path = C:/Windows/Fonts/,
    Extension = .ttf,
    BoldFont = segoeuib,
    ItalicFont = segoeuii,
    BoldItalicFont = segoeuiz]
  \setsansfont{segoeui}[
    Path = C:/Windows/Fonts/,
    Extension = .ttf,
    BoldFont = segoeuib,
    ItalicFont = segoeuii,
    BoldItalicFont = segoeuiz]
}{%
  \setmainfont{Latin Modern Sans}
  \setsansfont{Latin Modern Sans}
}

\usepackage[english,italian]{babel}
\AtBeginDocument{\selectlanguage{\ReportLang}}

\usepackage[a4paper,top=2.3cm,bottom=2.2cm,left=2.2cm,right=2.2cm,headheight=14pt,headsep=0.6cm,footskip=1cm]{geometry}
\usepackage[table]{xcolor}
\usepackage{colortbl}
\usepackage{array}
\usepackage{tabularx}
\usepackage{xltabular}
\usepackage{titlesec}
\usepackage{fancyhdr}
\usepackage{enumitem}
\usepackage{microtype}
\usepackage{needspace}
\usepackage{titletoc}
\usepackage[section]{placeins}
\usepackage{amsmath}
\usepackage{amssymb}
\usepackage{graphicx}
\usepackage{caption}
\usepackage{siunitx}

% Colori della nota metodologica
\definecolor{navy}{HTML}{0E2D53}
\definecolor{maroon}{HTML}{A3274B}
\definecolor{rowgrey}{HTML}{F2F2F2}
\definecolor{rulegrey}{HTML}{BFBFBF}
\definecolor{footgrey}{HTML}{8C8C8C}
\definecolor{boxgrey}{HTML}{F4F5F8}

\setlength{\parindent}{0pt}
\setlength{\parskip}{0.55em}
\linespread{1.12}

% Numeri: in italiano punto per le migliaia e virgola per i decimali, in inglese
% il contrario. Nel sorgente si scrive sempre il numero grezzo: \num{161250},
% \num{0.698}, \pct{96.6}. Il raggruppamento parte da 4 cifre (1.626).
% Modo testo: le cifre usano il font del testo e ereditano il grassetto.
% Incertezza esplicita: \num{0.874 +- 0.006} produce 0,874 ± 0,006.
\sisetup{mode=text,reset-text-family=false,reset-text-series=false,reset-text-shape=false,uncertainty-mode=separate}
\ifdefstring{\ReportLang}{english}{%
  \sisetup{output-decimal-marker={.}, group-separator={,}, group-minimum-digits=4}
  \newcommand{\ReportPageLabel}{p.}
}{%
  \sisetup{output-decimal-marker={,}, group-separator={.}, group-minimum-digits=4}
  \newcommand{\ReportPageLabel}{pag.}
}
\newcommand{\pct}[1]{\num{#1}\%}

% Testatina bordeaux in maiuscolo a sinistra, piede grigio a destra
\providecommand{\ReportHeader}{\ReportBrand\ifdefempty{\ReportProject}{}{\ \ | \ \ReportProject}}
\providecommand{\ReportFooter}{\ReportBrand\ifdefempty{\ReportProject}{}{ | \ReportProject}}
\newcommand{\ReportHeadStyle}{\footnotesize\bfseries\color{maroon}\MakeUppercase{\ReportHeader}}
\newcommand{\ReportFootStyle}{\scriptsize\color{footgrey}\ReportFooter\ ·\ \ReportPageLabel\ \thepage}
\pagestyle{fancy}
\fancyhf{}
\renewcommand{\headrulewidth}{0pt}
\fancyhead[L]{\ReportHeadStyle}
\fancyfoot[R]{\ReportFootStyle}
\fancypagestyle{plain}{\fancyhf{}\renewcommand{\headrulewidth}{0pt}%
  \fancyhead[L]{\ReportHeadStyle}\fancyfoot[R]{\ReportFootStyle}}

% Blocco del titolo: titolo blu in grassetto, sottotitolo, autore e ruolo, data
\newcommand{\maketitleblock}{%
  {\color{navy}\bfseries\LARGE \ReportTitle\par}
  \vspace{0.35em}
  {\color{navy}\large \ReportSubtitle\par}
  \vspace{0.5em}
  {\small \ReportAuthor, \ReportRole\par
  \color{footgrey}\ReportDate\par}}

% Titoli di sezione: numero e titolo blu in grassetto con filetto bordeaux sotto.
% Le sottosezioni sono blu, in grassetto, di corpo normale (piu' piccole delle sezioni).
\newcommand{\sectitle}[1]{\parbox[b]{\linewidth}{\thesection.\hspace{0.4em}#1\par\nointerlineskip\vspace{3pt}{\color{maroon}\rule{\linewidth}{0.7pt}}}}
\newcommand{\sectitlenn}[1]{\parbox[b]{\linewidth}{#1\par\nointerlineskip\vspace{3pt}{\color{maroon}\rule{\linewidth}{0.7pt}}}}
\titleformat{\section}[block]{\color{navy}\bfseries\large}{}{0pt}{\sectitle}
\titleformat{name=\section,numberless}[block]{\color{navy}\bfseries\large}{}{0pt}{\sectitlenn}
\titlespacing*{\section}{0pt}{1.0em}{0.2em}
\titleformat{\subsection}{\color{navy}\bfseries\normalsize}{\thesubsection}{0.6em}{}
\titlespacing*{\subsection}{0pt}{0.9em}{0.15em}
% Indice: sezioni in blu scuro, sottosezioni con puntini
\titlecontents{section}[1.6em]{\addvspace{0.4em}\color{navy}\bfseries}{\contentslabel{1.6em}}{\hspace*{-1.6em}}{\hfill\contentspage}
\titlecontents{subsection}[4.0em]{}{\contentslabel{2.4em}}{\hspace*{-2.4em}}{\titlerule*[0.7pc]{.}\contentspage}
% Evita titoli orfani in fondo pagina: se restano poche righe si va a capo pagina
\pretocmd{\section}{\needspace{9\baselineskip}}{}{}
\pretocmd{\subsection}{\needspace{5\baselineskip}}{}{}
% Le tabelle restano nella sezione in cui sono citate: placeins[section] chiude i
% float a ogni \section; a fine sezione si aggiunge comunque \FloatBarrier esplicito.
\renewcommand{\topfraction}{0.9}
\renewcommand{\bottomfraction}{0.8}
\renewcommand{\textfraction}{0.07}
\renewcommand{\floatpagefraction}{0.75}
\setcounter{topnumber}{3}
\setcounter{bottomnumber}{2}
\setcounter{totalnumber}{5}

\setlist{itemsep=0.15em,topsep=0.2em,parsep=0pt,leftmargin=1.3em}

% Didascalie delle tabelle: sopra la tabella, "Tabella N." in blu grassetto
\DeclareCaptionFont{navy}{\color{navy}}
\captionsetup{font=small,labelfont={bf,navy},textfont=navy,labelsep=period,justification=raggedright,singlelinecheck=false,skip=4pt,position=top}
\setlength{\LTcapwidth}{\linewidth}

% Tabelle: intestazione blu con testo bianco, righe alternate grigie, filetti sottili.
% Uso: \begin{tabella} ... tabularx o xltabular ... \end{tabella};
% riga di intestazione: \rowcolor{navy}\hd{Colonna} & ... \\ \hline
% xltabular va a capo pagina e ripete l'intestazione posta prima di \endhead.
\arrayrulecolor{rulegrey}
\setlength{\arrayrulewidth}{0.4pt}
\setlength{\LTpre}{0.3em}
\setlength{\LTpost}{0.3em}
\newcommand{\hd}[1]{\textcolor{white}{\bfseries #1}}
\newcolumntype{L}{>{\raggedright\arraybackslash}X}
\newcolumntype{P}[1]{>{\raggedright\arraybackslash}p{#1}}
\newcolumntype{R}[1]{>{\raggedleft\arraybackslash}p{#1}}
\newcolumntype{C}[1]{>{\centering\arraybackslash}p{#1}}
\newenvironment{tabella}{\par\footnotesize\renewcommand{\arraystretch}{1.3}\rowcolors{2}{rowgrey}{white}}{\par\normalsize}

% Riquadro di sintesi: bordo blu, fondo grigio chiaro (non si spezza tra pagine)
\newsavebox{\riquadrobox}
\newenvironment{riquadro}
  {\par\noindent\begin{lrbox}{\riquadrobox}\begin{minipage}{\dimexpr\linewidth-2\fboxsep-2\fboxrule\relax}\vspace{0.3em}}
  {\vspace{0.2em}\end{minipage}\end{lrbox}\fcolorbox{navy}{boxgrey}{\usebox{\riquadrobox}}\par}

% Collegamenti: hyperref per ultimo
\usepackage{hyperref}
\hypersetup{colorlinks=true,linkcolor=navy,citecolor=navy,urlcolor=maroon}
\AtBeginDocument{\hypersetup{pdftitle={\ReportTitle},pdfauthor={\ReportAuthor}}}
% cleveref dopo hyperref: \cref{tab:x} produce "Tabella 1"
\usepackage[capitalise,noabbrev]{cleveref}
